\section{Introduction}
\label{sec:introduction}

The Efficiency Gap became the dominant partisan fairness metric in American redistricting
litigation between its introduction by \citet{stephanopoulosMcGhee2015} and the Supreme
Court's foreclosure of federal partisan gerrymandering claims in \textit{Rucho v.\
Common Cause}.\footnote{\textit{Rucho v.\ Common Cause}, 588 U.S.\ 684 (2019).}
Its central appeal is the intuitive foundation: democratic elections should not
systematically waste more votes from one party than the other, and the EG
measures exactly this asymmetry in a single number that is computable from
publicly available election returns.

EG rose to prominence in \textit{Gill v.\ Whitford},\footnote{\textit{Gill v.\
Whitford}, 585 U.S.\ 48 (2018).} where plaintiffs used it as the primary
quantitative evidence that Wisconsin's 2012 state assembly maps were a partisan
gerrymander. The district court ruled for the plaintiffs partly on EG evidence;
SCOTUS reversed on standing grounds, declining to reach the merits. \textit{Rucho}
then foreclosed the merits question entirely at the federal level. This sequence
left EG in an unusual legal position: extensively discussed in federal judicial
opinions, endorsed by leading political scientists, proposed as a constitutional
standard---but ultimately not adopted by any federal court.

State courts have continued to cite EG in partisan gerrymandering claims, but
without any binding threshold. The North Carolina Supreme Court in \textit{Harper
v.\ Hall}\footnote{\textit{Harper v.\ Hall}, 868 S.E.2d 499 (N.C.\ 2022).}
treated EG as one of several indicators of partisan manipulation, not as a
standalone standard. The result is that EG remains a central piece of expert
testimony evidence---widely computed, widely cited, and well understood by
redistricting courts---without any legal framework specifying what value makes
a map constitutionally infirm.

This paper fills the need for a rigorous analysis of EG's behavior under
\textsc{bisect} algorithmic redistricting. The \textsc{bisect} algorithm,
which draws congressional districts using only population counts and geographic
adjacency without any partisan data, produces maps with EG values that provide
a natural reference point for state court comparison. Section~\ref{sec:definition}
develops the mathematical definition and its properties. Section~\ref{sec:behavior}
explains why compactness-driven redistricting systematically reduces EG.
Section~\ref{sec:empirical} presents multi-state, multi-election empirical
results with sensitivity analysis. Section~\ref{sec:legal} provides a
comprehensive legal history from the metric's introduction through 2026.
Section~\ref{sec:conclusion} concludes with guidance for expert witnesses.
